% sec-02-introduction.tex - Introduction (full)
\section{Introduction}

\subsection{Introductory Statement}

This Investigational New Drug (IND) application (IND v1.0, repository v4.3.0)
supports a Phase 1, first-in-human, combined IND and Investigational Device
Exemption (IDE) clinical investigation of on-premises large language model (LLM)
directed eight-arm robotic pancreaticoduodenectomy (Whipple) with perioperative
daraxonrasib (RMC-6236) in patients with KRAS G12-mutated pancreatic ductal
adenocarcinoma (PDAC). The investigation is filed under 21 CFR \S 312 for the
investigational drug and is coupled to a significant-risk early-feasibility device
study under 21 CFR \S 812 for the LLM-directed robotic platform; the two pathways
are coordinated within a single integrated protocol because the device conducts the
resection that generates the operative and pathology inputs to the daraxonrasib
restart advisory \cite{daraxwhipple, onpremwhippl}. The sponsor-investigator is
ChemicalQDevice (Kevin Kawchak, Chief Executive Officer), and the conduct,
oversight, and reporting commitments described throughout this dossier are made in
accordance with the applicable provisions of 21 CFR \S 312.23 and the FDA Phase 1
content and format guidance \cite{phase1guidance, cfr312paiuni}.

The clinical problem this investigation addresses is severe and time-sensitive.
PDAC is the third leading cause of cancer death in the United States and the fourth
in the European Union, with total five-year overall survival below 13 percent
\cite{pdac060s2030}. The pancreaticoduodenectomy is the most technically demanding
of all curative-intent abdominal oncology operations: it requires resection and
reconstruction of four luminal structures (the pancreas, the duodenum, the common
bile duct, and the stomach or pylorus) and dissection around five named vessels
(the superior mesenteric vein, the portal vein, the hepatic artery, the celiac
axis, and the superior mesenteric artery). The leading lethal postoperative
sequence is the fistula-sepsis cascade: a clinically relevant pancreatic fistula
that progresses to intra-abdominal infection, hemorrhage, sepsis, and failure to
rescue. The 2025 Dutch nationwide cohort of 1000 robotic
pancreaticoduodenectomies, which anchors the contemporary human baseline, reported
a conversion rate of 10.1 percent, Clavien-Dindo grade III or higher complications
of 41.3 percent, International Study Group on Pancreatic Surgery (ISGPS) grade B or
C fistula of 24.4 percent, and 30-day mortality of 3.9 percent. Even expert,
high-volume robotic surgery therefore leaves a substantial residual of conversions,
high-grade complications, and clinically relevant fistulae. The intervention
studied here pairs a precise, auditable, on-premises LLM-directed robotic Whipple
with the RAS(ON) multi-selective inhibitor daraxonrasib, targeting an in-window R0
resection while layering hard cyber-physical safety limits and optimizing the
perioperative drug-restart window.

The device that conducts the resection is an eight-arm parallel coelomic surgical
platform carrying 56 mechanical degrees of freedom (eight arms with seven
degrees of freedom each) and a 640-channel mixed-rate sensor stack (80 channels per
arm), with force channels sampling at 100 kHz and all remaining channels at the
10 kHz command rate, positioning accuracy of 0.05 mm root mean square at a peak tip
velocity of 1200 mm/s. An on-premises repository-pinned LLM acts as a second-opinion
advisory oracle held strictly outside the robot vendor kinematic stack and hash-pinned
to a deterministic seed and a specific commit, never as an autonomous controller;
full autonomy is prohibited consistent with 21 CFR \S 312.21(e). The hard
cyber-physical safety envelope comprises a per-arm tip-force cap of 3 N or less, a
cumulative cross-arm tip-force cap of 18 N or less enforced at the heartbeat
boundary, a 10 kHz heartbeat coordination bus broadcasting a 64-byte frame on a 100
microsecond per-arm watchdog deadline, an emergency arm park within 50 microseconds
of a missed or out-of-parameter frame, a cross-arm emergency stop that halts the
full platform in 3 ms or less, and a software-independent hardware emergency stop in
500 ms or less consistent with 21 CFR \S 312.404. A vascular no-fly gate samples at
10 kHz around five named vessels (the superior mesenteric vein, portal vein, hepatic
artery, celiac axis, and superior mesenteric artery), with a soft-warning advisory
zone at 3.0 mm, a no-fly blocking zone at 1.0 mm for the superior mesenteric and
portal veins and 1.5 mm for the hepatic artery, celiac axis, and superior mesenteric
artery, and a hard-stop emergency-stop zone at 5.0 mm. In this Phase 1 the device
operates only at the clinician-controlled and supervised-semiautonomous levels under
continuous human oversight with immediate manual override, consistent with the
staged-autonomy posture in which higher autonomy is reserved for future stages
\cite{daraxwhipple, onpremwhippl, clintrustpai}.

The eligible population is adults (18 years or older) with histologically or
cytologically confirmed, resectable or borderline-resectable PDAC; documented KRAS
G12 mutation (G12D, G12V, or G12R); ECOG performance status 0 or 1; and candidacy
for pancreaticoduodenectomy, eligible under the broad pan-KRAS criteria of the
RASolute 302 program. An exemplar screening subject, PAT-PDAC-0001, is a 68-year-old
with a 3.4 cm head-of-pancreas tumor abutting the superior mesenteric vein at 75
degrees, KRAS G12D, microsatellite-stable disease, CA 19-9 of 412 U/mL, ECOG 1, and
four cycles of neoadjuvant modified FOLFIRINOX, illustrating the borderline-resectable
phenotype the investigation targets.

The investigational plan rests on a three-part acceleration thesis that is the
animating premise of the present repository. A patient already enrolled in a Phase
1 PDAC study has little survival margin to spare on administrative delay; the
clinical and operational time of a trial (recruitment, treatment, and waiting for
progression-free and overall survival to mature) cannot be compressed by writing
faster, and neither can the fixed regulatory review clocks, but the administrative
and preparation time spent authoring the large documents between steps can be
compressed substantially \cite{oncotrialllm, paipredict4x}. Faster, internally
consistent paperwork therefore means faster treatment: a complete and consistent
package reaches the Institutional Review Board (IRB) and the FDA sooner, the 30-day
IND safety clock starts sooner, sites activate sooner, and the next dose-escalation
cohort opens sooner. This IND, like the protocol and the supporting publication
from which it is assembled, was generated through a single-prompt mermaid to draft
to full to final pipeline whose every artifact is committed and reviewable in real
time, and the same mechanism that produced this dossier maintains the safety
reports, amendments, and study reports that the trial will require throughout its
lifecycle.

Figure 1 places this IND within the full Phase 1 oncology document landscape it
inhabits. Before submission the repository LLM generates the IND dossier, the
Investigator's Brochure, the clinical protocol, and the informed consent form;
during the trial it produces the 7-day and 15-day safety reports, the annual
Development Safety Update Report (DSUR), protocol amendments, and dose-escalation
minutes; and after the trial it assembles the tables, listings, and figures, the
ICH E3 clinical study report, and the downstream manuscripts and lay summaries. The
IND package is the first and rate-limiting artifact in this relay, which is why the
acceleration thesis is most consequential here \cite{oncotrialllm}.

\begin{mermaidfig}
\node[mmaccent](hub) at (5.4,-3.0)  {Repository LLM\\IND document\\generation};
\node[mmin]    (b1)  at (0,0)       {IND dossier\\(1571 to 11)};
\node[mmin]    (b2)  at (0,-1.9)    {Investigator's\\Brochure};
\node[mmin]    (b3)  at (0,-3.8)    {Clinical protocol};
\node[mmin]    (b4)  at (0,-5.7)    {Informed consent\\(ICF)};
\node[mmproc]  (d1)  at (10.8,0)    {SAE / SUSAR\\7 / 15-day};
\node[mmproc]  (d2)  at (10.8,-1.9) {DSUR (annual)};
\node[mmproc]  (d3)  at (10.8,-3.8) {Protocol\\amendments};
\node[mmproc]  (d4)  at (10.8,-5.7) {Dose-escalation\\minutes};
\node[mmgoal]  (a1)  at (5.4,-7.4)  {Tables, listings,\\figures};
\node[mmgoal]  (a2)  at (5.4,-9.3)  {CSR (ICH E3)};
\node[mmgoal]  (a3)  at (5.4,-11.2) {Manuscripts /\\lay summaries};
\draw[mmedge] (b1.east) to[out=0,in=160,looseness=0.6] (hub.north west);
\draw[mmedge] (b2.east) to[out=0,in=175,looseness=0.6] (hub.west);
\draw[mmedge] (b3.east) to[out=0,in=185,looseness=0.6] (hub.west);
\draw[mmedge] (b4.east) to[out=0,in=200,looseness=0.6] (hub.south west);
\draw[mmedge] (hub.north east) to[out=20,in=180,looseness=0.6] (d1.west);
\draw[mmedge] (hub.east) to[out=5,in=180,looseness=0.6] (d2.west);
\draw[mmedge] (hub.east) to[out=-5,in=180,looseness=0.6] (d3.west);
\draw[mmedge] (hub.south east) to[out=-20,in=180,looseness=0.6] (d4.west);
\draw[mmedge] (hub) -- (a1);
\draw[mmedge] (a1) -- (a2);
\draw[mmedge] (a2) -- (a3);
\end{mermaidfig}
\vspace{-0.3cm}
\figcaption{Figure 1. The large Phase 1 oncology document landscape around the IND,
grouped into before-submission inputs (the IND dossier, Investigator's Brochure,
protocol, and consent), during-trial process documents (7 and 15-day safety
reports, the annual DSUR, amendments, and dose-escalation minutes), and after-trial
outputs (tables, listings, figures, the ICH E3 clinical study report, and
manuscripts and lay summaries), all generated by the repository LLM around the
central hub. This figure renders in the Introduction to place the IND in the full
document lifecycle the same method maintains.}

The remainder of this Introduction follows the ReGARDD IND content sequence. It
names the drug and its active ingredients (\S 3.1.1), states the pharmacological
class (\S 3.1.2) and structural description (\S 3.1.3), summarizes the dosage form
(\S 3.1.4), the route of administration (\S 3.1.5), and the duration of the
proposed clinical investigation (\S 3.1.6); it then summarizes previous human
experience (\S 3.2), the status of the drug in other countries (\S 3.3), and the
preclinical and computational evidence (\S 3.4), with the supporting references
collected in \S 3.5. Quantitative and design detail introduced here is developed
in full in the General Investigational Plan (\S 4), the Investigator's Brochure
(\S 5), the proposed clinical research (\S 6), the Chemistry, Manufacturing, and
Control information (\S 7), and the Pharmacology and Toxicology information (\S 8).

\subsubsection{Name of the Drug and All Active Ingredients}

The investigational drug is daraxonrasib, also known by its development code
RMC-6236. Daraxonrasib is the single active ingredient of the drug product; there
are no additional active pharmaceutical ingredients. It is an orally bioavailable,
RAS(ON) multi-selective (pan-RAS) small-molecule inhibitor that binds the active,
guanosine-triphosphate-bound (ON) state of mutant RAS in a tri-complex with the
intracellular chaperone cyclophilin A, suppressing downstream effector signaling
across the common KRAS oncoprotein variants found in PDAC \cite{daraxwhipple,
onpremwhippl}. The molecular target population for this investigation is defined by
the KRAS G12 alleles (G12D, G12V, and G12R) that predominate in pancreatic ductal
adenocarcinoma; daraxonrasib provides broad pan-KRAS coverage that spans these
alleles, which is the property that makes it suitable for the protocol's
molecularly defined eligibility \cite{daraxwhipple}.

Figure 2 reproduces the mechanism and pharmacological-class schematic for
daraxonrasib that grounds \S 3.1.1 through \S 3.1.3. Daraxonrasib forms the
tri-complex with cyclophilin A, binds the active GTP-bound state of mutant RAS,
intercepts the KRAS G12 variants that define the trial population, suppresses
downstream MAPK and PI3K effector signaling, and thereby reduces tumor-cell
proliferation and survival; it is dosed orally once daily at the three escalation
levels (DL1 160 mg, DL2 220 mg, DL3 300 mg) and carries FDA Breakthrough Therapy
designation granted in June 2025 for previously treated metastatic KRAS G12 PDAC
\cite{daraxwhipple, onpremwhippl}.

\begin{mermaidfig}
\node[mmgoal] (drx)   at (0,0)       {Daraxonrasib (RMC-6236)\\oral RAS(ON)\\multi-selective\\(pan-RAS) inhibitor};
\node[mmproc] (cypa)  at (0,-2.6)    {Tri-complex with\\cyclophilin A};
\node[mmproc] (rason) at (4.6,-2.6)  {Binds active GTP-bound\\mutant RAS (RAS(ON))};
\node[mmin]   (kras)  at (4.6,0)     {KRAS G12 variants\\G12D / G12V / G12R};
\node[mmaccent](eff)  at (4.6,-5.2)  {Suppresses downstream\\MAPK / PI3K effector\\signaling};
\node[mmgoal] (prol)  at (0,-5.2)    {Reduced tumor-cell\\proliferation and survival};
\node[mmctx]  (route) at (9.4,-1.3)  {Route: oral, once daily\\DL1 160 / DL2 220 /\\DL3 300 mg};
\node[mmdark] (btd)   at (9.4,-4.1)  {FDA Breakthrough Therapy\\(June 2025), pretreated\\metastatic KRAS G12 PDAC};
\draw[mmedge] (drx) -- (cypa);
\draw[mmedge] (cypa) -- (rason);
\draw[mmedge] (kras) -- (rason);
\draw[mmedge] (rason) -- (eff);
\draw[mmedge] (eff) -- (prol);
\draw[mmedged] (drx.east) to[out=0,in=160,looseness=0.7] (route.west);
\draw[mmedged] (drx.south east) to[out=-30,in=170,looseness=0.6] (btd.west);
\end{mermaidfig}
\vspace{-0.3cm}
\figcaption{Figure 2. Daraxonrasib (RMC-6236) is an oral RAS(ON) multi-selective
(pan-RAS) small-molecule inhibitor. It forms a tri-complex with cyclophilin A and
binds the active, guanosine-triphosphate-bound state of mutant RAS, including the
KRAS G12 variants (G12D, G12V, G12R) that define the trial population, suppressing
downstream MAPK and PI3K effector signaling and reducing tumor-cell proliferation
and survival. It is dosed orally once daily at DL1 160 mg, DL2 220 mg, and DL3 300
mg, and holds FDA Breakthrough Therapy designation (June 2025) for previously
treated metastatic KRAS G12 PDAC. This figure renders in \S 3.1.1 (Name of Drug and
Active Ingredients), \S 3.1.2 (Pharmacological Class), and \S 3.1.3 (Structural
Formula).}

In the present investigation daraxonrasib is studied in a perioperative,
curative-intent surgical context that is distinct from its late-phase metastatic
program: it is administered orally once daily in a defined preoperative course,
paused across the operative window, and restarted postoperatively on a day chosen
by the LLM-bound, human-confirmed advisory described in \S 3.1.5 and developed in
the proposed clinical research. Because the drug is paired with an investigational
device in the drug-delivery pathway, the combined IND and IDE submission carries a
full Physical AI System Description meeting the eleven fields of 21 CFR
\S 312.23(g) \cite{cfr312paiuni}, the substance of which is presented in the CMC
and device sections of this dossier.

\subsubsection{Pharmacological Class of the Drug}

Daraxonrasib belongs to the RAS(ON) multi-selective small-molecule inhibitor class.
Pharmacologically, the agent is a tri-complex inhibitor: it does not act on the
inactive, guanosine-diphosphate-bound (OFF) conformation that earlier-generation
allele-specific inhibitors target, but instead engages the active, GTP-bound (ON)
conformation of mutant RAS by recruiting the abundant intracellular chaperone
cyclophilin A to form a high-affinity ternary complex on the RAS surface
\cite{daraxwhipple, onpremwhippl}. This ON-state, chaperone-mediated mechanism
confers multi-selective (pan-RAS) coverage across KRAS, NRAS, and HRAS variants,
and in particular across the KRAS G12 alleles (G12D, G12V, G12R) that account for
the large majority of RAS-driven PDAC. By occupying the active oncoprotein, the
inhibitor blocks engagement of downstream effectors and suppresses the mitogen
activated protein kinase (MAPK) and phosphoinositide 3-kinase (PI3K) signaling
axes, with the consequence of reduced tumor-cell proliferation and survival, as
summarized in Figure 2.

The class assignment is material to the safety profile anticipated in this Phase 1
study and to the dose-finding rationale. As a multi-selective RAS pathway
inhibitor, daraxonrasib carries the on-target, mechanism-based toxicities
characteristic of broad RAS and MAPK pathway suppression, including
gastrointestinal effects (nausea, diarrhea, stomatitis), dermatologic and
mucocutaneous events (rash), and hepatic enzyme elevations; these are the toxicity
categories that the 3+3 escalation is specifically designed to detect and bound and
that inform the dose-limiting toxicity (DLT) definitions in \S 4. The class also
informs the perioperative pause-and-restart logic, because broad pathway inhibition
during the immediate postoperative window carries a theoretical risk to wound and
anastomotic healing if the restart is mistimed, which is why the restart day is
keyed to fistula grade and drug trough rather than fixed.

\subsubsection{Structural Formula of the Drug}

Daraxonrasib (RMC-6236) is a synthetic, non-peptidic small-molecule organic
compound of defined molecular structure designed to bind the RAS-cyclophilin A
interface and engage the active GTP-bound state of mutant RAS. It is a single,
well-characterized chemical entity rather than a mixture, a biologic, or a
nucleic-acid therapeutic, and it is the sole active ingredient of the drug product.
The molecule is the macrocyclic tri-complex inhibitor scaffold characteristic of
the RAS(ON) class, engineered for oral bioavailability and once-daily exposure, and
its structure supports the broad pan-KRAS coverage that defines the molecularly
selected trial population.

The complete structural formula of daraxonrasib, including the molecular formula,
molecular weight, stereochemistry, the structural diagram of the macrocyclic
scaffold, the chemical name, and the characterization data establishing identity
and purity, is provided in the Chemistry, Manufacturing, and Control information of
this submission at \S 7.1.1 (Drug Substance), consistent with 21 CFR \S 312.23(a)(7)
\cite{phase1guidance, cfr312paiuni}. Presenting the full structural characterization
in the CMC section avoids duplication and keeps the structural, physicochemical, and
manufacturing detail consolidated where the reviewer expects it; this Introduction
provides only the qualitative structural description and the cross-reference. The
mechanism schematic in Figure 2 summarizes the functional relationship between the
structure and its pharmacological target.

\subsubsection{Formulation of the Dosage Form(s) to be Used}

Daraxonrasib is supplied as oral tablets in the strengths required to compose the
three once-daily dose levels evaluated in the 3+3 escalation: DL1 at 160 mg, DL2 at
220 mg, and DL3 at 300 mg. The 300 mg once-daily level (DL3) corresponds to the
recommended Phase 2 dose established in the pan-KRAS RASolute 302 program, and the
escalation begins one and two steps below it (160 mg and 220 mg) to preserve a
conservative first-in-human margin for the novel perioperative surgical context
\cite{daraxwhipple, onpremwhippl}. Tablets are swallowed whole with water and are
not crushed, split, or chewed unless a future protocol amendment provides validated
instructions. No compounding or reconstitution is required; preparation consists of
the site pharmacist verifying participant identity, the assigned dose level, and the
expiry, then dispensing the once-daily oral dose with administration instructions.

Tablets are packaged in child-resistant, light-protective containers labeled in
accordance with 21 CFR \S 312.6, bearing the investigational caution statement
``Caution: New Drug. Limited by Federal (or United States) law to investigational
use,'' together with the protocol number, lot number, storage conditions, and a
unique container identifier; because the study is open-label, the participant-facing
label identifies the product and the assigned daily dose. The product is stored at
controlled room temperature in a secured, access-controlled, temperature-monitored
location within the investigational-product pharmacy, in its original
light-protective packaging, with continuous temperature monitoring and documented
excursion handling, and is dispensed only within its labeled expiry or retest date.
The dose-level structure that the dosage form must support is summarized in Table 1;
the full formulation composition, appearance, packaging, stability, and labeling
detail is provided in \S 7.1.2 (Drug Product) and \S 7.1.4 (Labeling).

\begin{table}[t]
\caption{Daraxonrasib dose-level structure for the Phase 1 3+3 escalation. DL3 (300
mg once daily) corresponds to the pan-KRAS RASolute 302 recommended Phase 2 dose;
DL1 and DL2 sit below it to preserve a conservative first-in-human margin. The DLT
window is 28 days, and the recommended Phase 2 dose (RP2D) is selected at or below
the maximum tolerated dose (MTD).}
\label{tab:sec02-doselevels}
\centering
\begin{tabularx}{\textwidth}{L{1.6cm} L{2.5cm} C{1.7cm} L{3.0cm} Y}
\toprule
\textbf{Dose level} & \textbf{Daily dose} & \textbf{DLT window} & \textbf{3+3 cohort} & \textbf{Role and rationale} \\
\midrule
DL1 & 160 mg QD & 28 days & 3, expand to 6 if 1 DLT & Conservative starting dose, two steps below the RASolute 302 reference \\
DL2 & 220 mg QD & 28 days & 3, expand to 6 if 1 DLT & Intermediate level, one step below the reference dose \\
DL3 & 300 mg QD & 28 days & 3, expand to 6 if 1 DLT & RASolute 302 recommended Phase 2 dose; highest planned level \\
\bottomrule
\end{tabularx}
\end{table}

\figcaption{Table 1. Daraxonrasib oral dose levels and the 3+3 escalation
structure. At each level three participants are observed across the 28-day DLT
window; the cohort escalates on 0 of 3, expands to 6 on 1 of 3, escalates on 1 of
6, and fixes the MTD at the preceding level on 2 or more of 6.}

\subsubsection{Route of Administration}

Daraxonrasib is administered by the oral route, once daily, on a perioperative
schedule comprising a defined preoperative course, a protocol-mandated pause
spanning the operative window, and a postoperative restart whose timing is governed
by an advisory. The preoperative course establishes systemic exposure before
surgery; the pause spans the immediate operative and early-recovery window during
which broad RAS pathway suppression could impair wound and anastomotic healing; and
the restart returns systemic therapy at the earliest safe point to suppress
micrometastatic outgrowth. The drug is not delivered by the device, which has no
pharmacologic dose of its own; the device conducts the resection and reconstruction
and supplies the operative and pathology inputs that the restart advisory consumes
\cite{daraxwhipple}.

The perioperative pause-and-restart advisory is LLM-bound and human-confirmed, never
autonomous. It takes two inputs and produces one of three restart recommendations.
The inputs are the pancreaticojejunostomy ISGPS postoperative pancreatic-fistula
grade (A, B, or C) and the serum daraxonrasib trough relative to a 0.5 ng/mL
threshold. The outputs are a restart on postoperative day T+7d, T+14d, or T+21d, or
a continued hold. In the prespecified 32-iteration deterministic simulation sweep,
baseline serum was 0.45 ng/mL at the operative timepoint across all iterations
(below the 0.5 ng/mL trough), and the advisory recommended T+7d in 29 of 32
iterations (90.6 percent, Grade A fistula with sub-threshold trough), T+14d in 3 of
32 iterations (9.4 percent, Grade B fistula or a delayed serum rise), and T+21d in 0
of 32 iterations; Grade C fistula maps to T+21d or a continued hold and to the drug
stop rule. Because the pancreaticojejunostomy is the dominant determinant of
postoperative pancreatic-fistula risk, the fistula-grade input to the drug advisory
is based on that anastomosis. The advisory is one of only two coupling points
between the drug and device components of the integrated intervention, and it keeps
the restart decision human-confirmed and fully traceable in the audit trail. The
detailed advisory logic and the simulation sweep are presented in the proposed
clinical research and Investigator's Brochure sections of this dossier.

\subsubsection{Duration of the Proposed Clinical Investigation(s)}

The duration of the proposed clinical investigation is defined at two levels: the
per-participant timeline and the overall study timeline. For each participant, the
investigation comprises up to 28 days of screening (Day -28 to Day -1) during which
histologic confirmation, KRAS G12 confirmation, staging, and eligibility are
established; a baseline day (Day -1) that carries the Phase 0 simulation sign-off
and the Unified Safety Level (USL) and safety-matrix verification; surgery on Day 0;
an acute window (Day 1 to Day 7) carrying the ISGPS fistula grading and the
daraxonrasib restart advisory; the Day 30 primary safety window during which
device-related and procedure-related serious adverse events, including Clavien-Dindo
grade III or higher complications, are ascertained; Day 90 pathology and mortality
assessment, which carries the R0 resection determination and the 90-day mortality
endpoint; and long-term follow-up every 12 weeks to the 24-month overall-survival
endpoint, during which progression-free survival (PFS) and overall survival (OS)
are assessed by RECIST v1.1 \cite{daraxwhipple}.

The daraxonrasib DLT observation window is 28 days from the start of dosing at each
dose level, which governs the cadence of the 3+3 escalation. At the study level, the
investigation enrolls up to 18 treated participants (approximately 36 screened) at a
single academic site, with an anticipated accrual of approximately one to two
participants per month under sentinel staggered enrollment; with three dose levels
of up to six DLT-evaluable participants each, the maximum of 18 follows directly
(three levels times six). Adding the per-participant follow-up to 24-month OS to the
accrual period, the total study duration is governed primarily by the time required
to mature the overall-survival endpoint in the last-enrolled participant. A
mandatory Phase 0 simulation-validation gate (at least 1000 simulated procedures
across at least two independent frameworks, with a USL rating of at least 7.0 and a
sim-to-real trajectory gap below 2 mm and tip-force gap below 0.5 N) precedes any
patient-facing robotic activity and is re-validated after any qualifying change, so
the patient-facing portion of the investigation does not begin until that gate is
satisfied. The first-year operational plan and the projected enrollment cadence are
developed in the General Investigational Plan (\S 4).

The investigation is statistically descriptive and non-confirmatory, consistent with
ICH E9 principles for early-phase trials and with TRIPOD+AI for the AI-enabled
components; there is no power calculation and no comparator arm. Estimation uses
exact Clopper-Pearson 95 percent confidence intervals, and the maximum sample of 18
participants yields the precision appropriate to a first-in-human dose-finding and
early-feasibility study: an observed proportion of 0 of 18 corresponds to 0 percent
(95 percent CI 0 to 18 percent), 1 of 18 corresponds to 5.6 percent (95 percent CI
0.1 to 27 percent), 0 of 3 to a CI of 0 to 71 percent, and 0 of 6 to a CI of 0 to 46
percent. The three co-primary endpoints are the incidence of device-related and
procedure-related serious adverse events through the Day 30 window (including
Clavien-Dindo grade III or higher), the maximum tolerated dose and recommended Phase
2 dose of daraxonrasib determined by the 3+3 escalation across the 28-day DLT window,
and technical feasibility defined as the proportion of participants completing the
assigned operative tasks without unplanned conversion attributable to device
malfunction or unsafe behavior. Binding halt rules govern progression: a
device-related or procedure-related serious adverse event in at least one of three
participants within a cohort, or two device-related deaths at any time, triggers Data
Safety Monitoring Board review and possible halt. Safety reporting follows the
7-calendar-day clock for fatal or life-threatening suspected adverse reactions and
the 15-calendar-day clock for other serious and unexpected reactions under 21 CFR
\S 312.32, augmented by six Physical AI reporting triggers under 21 CFR \S 312.32(g)
and a -24 hour to +72 hour audit-trail preservation window around each Physical AI
adverse event \cite{phase1guidance, cfr312paiuni}.

\subsection{Summary of Previous Human Experience}

Daraxonrasib has an active and advancing clinical development program in
RAS-mutated solid tumors, with the most mature human experience in pancreatic
cancer. The FDA granted daraxonrasib Breakthrough Therapy designation in June 2025
for previously treated metastatic PDAC with KRAS G12 mutations, reflecting
preliminary clinical evidence of a meaningful advantage over available therapy in
this population \cite{daraxwhipple, onpremwhippl}. The late-phase program comprises
RASolute 302, a Phase 3 study in previously treated metastatic PDAC, and RASolve
301, a first-line program; the broad pan-KRAS eligibility of RASolute 302 defines
the molecular eligibility adopted in this protocol, and 300 mg once daily is the
established recommended Phase 2 dose carried forward as the highest dose level (DL3)
in the present escalation \cite{daraxwhipple}. The previous human experience with
daraxonrasib therefore establishes a pan-KRAS antitumor rationale, a tolerability
and exposure basis for the dose levels studied here, and a regulatory recognition
(Breakthrough Therapy) of unmet need in the same molecular population. The
perioperative pause-and-restart logic additionally respects any RASolve 301
checkpoint-inhibitor combination context.

The present investigation differs from the prior daraxonrasib human experience in
two material respects that define the scope of the first-in-human determination for
this IND. First, it studies the drug in a perioperative, curative-intent surgical
setting (resectable and borderline-resectable disease) rather than in the
metastatic, palliative setting of the late-phase program, so the wound-healing and
anastomotic-healing considerations are new and motivate the pause-and-restart
advisory. Second, it couples the drug to an investigational LLM-directed robotic
device, which has no prior human experience and which is evaluated here as a
significant-risk early-feasibility Class II collaborative platform under the IDE
pathway; the device-related and procedure-related safety endpoints in \S 4 are
specific to this novel combination. The full marketed-experience, prior-clinical,
and clinical-care experience analysis is presented in the Previous Human Experience
section of this dossier (\S 9).

The investigational design is further grounded by a body of computational and
in silico evidence developed by the sponsor-investigator throughout 2025 and 2026,
which supplies the simulation scaffolding adapted in this protocol and independently
complements the daraxonrasib drug-development program. This evidence comprises a
60-second PDAC robotic Whipple and daraxonrasib simulation; an AI digital-twin
pancreatic-cancer simulation for in silico FDA compliance and cost efficiency
\cite{fdadigtwinpc}; a quantitative systems pharmacology (QSP) metastatic-PDAC trial
simulation with code, verification, validation, and uncertainty quantification, and
a reproducible playbook \cite{qspmetpancre}; a 100,000-patient in silico Phase 3
five-arm PDAC trial triplicate \cite{chatgpt100kp}; and end-to-end PDAC digital-twin
clinical-trial proposals \cite{pdacdigtwinp}. The sponsor-investigator's record of
LLM evaluation and trust over August 2024 to June 2026 further grounds the device
governance, including a comparative code-generation evaluation of 16 proprietary
versus open-source LLMs against the FDA Adverse Event Reporting System
\cite{llmcomp16fda}, the on-premises advisory-LLM and verification-before-generation
design \cite{clintrustpai}, and the legislative framing aligned with H.R. 9510, the
Verification Before Generation in Physical AI Oncology Trials Act of 2026
\cite{congress9510}. The previous human experience and the supporting computational
evidence together motivate the favorable benefit-risk determination for this
low-survival, limited-alternative population.

\subsection{Status of Drug in Other Countries}

Daraxonrasib is an investigational agent that has not received marketing approval
in the United States or in any other jurisdiction for any indication; it is studied
worldwide under investigational authorizations within its sponsor's RAS(ON)
development program, and its U.S. Breakthrough Therapy designation (June 2025) is a
development-stage designation rather than a marketing authorization
\cite{daraxwhipple, onpremwhippl}. The drug has therefore not been withdrawn from
marketing in any country for reasons of safety or effectiveness, and there is no
adverse foreign regulatory action to report under 21 CFR \S 312.23(a)(5). Within the
present investigation, the LLM-directed robotic device component is similarly
investigational and is studied under the IDE pathway as a significant-risk Class II
collaborative device; it carries no marketing clearance or approval in any
jurisdiction.

Although this Phase 1 investigation is conducted at a single United States academic
site under the combined IND and IDE, the dossier is authored to align with the
principal international device and trial frameworks so that the same documentation
can support future multi-jurisdiction extension without re-derivation. The
cross-jurisdiction alignment that the device safety architecture and the
Physical AI System Description are designed to satisfy is summarized in Table 2. The
hard cyber-physical safety envelope (the 10 kHz heartbeat with a 100 microsecond
watchdog, the cross-arm emergency stop of 3 ms or less, the per-arm tip-force cap of
3 N or less and cumulative cap of 18 N or less, and the vascular no-fly gating) and
the 21 CFR part 11 hash-chained audit trail are constructed to map onto the
corresponding requirements of the European Union Medical Device Regulation, the
China National Medical Products Administration, the Japan Pharmaceuticals and
Medical Devices Agency, the United Kingdom Medicines and Healthcare products
Regulatory Agency, Health Canada, and the Australian Therapeutic Goods
Administration \cite{natlpaiplatf, cfr312paiuni}.

\begin{table}[t]
\caption{Cross-jurisdiction regulatory status and alignment. Daraxonrasib is
investigational with no marketing authorization in any country; the device is
investigational under the IDE. The dossier is authored to align with the principal
international frameworks for future multi-jurisdiction extension.}
\label{tab:sec02-jurisdiction}
\centering
\begin{tabularx}{\textwidth}{L{2.6cm} L{3.4cm} Y}
\toprule
\textbf{Jurisdiction} & \textbf{Framework} & \textbf{Status and alignment for this investigation} \\
\midrule
United States & FDA, 21 CFR \S 312 / \S 812 & Combined IND / IDE; daraxonrasib investigational with Breakthrough Therapy designation (June 2025); device significant-risk Class II under IDE \\
European Union & EU MDR 2017/745 & No marketing authorization; safety envelope and audit trail authored to map onto MDR essential requirements \\
China & NMPA & No marketing authorization; device safety architecture aligned for future NMPA submission \\
Japan & PMDA & No marketing authorization; Physical AI System Description aligned to PMDA expectations \\
United Kingdom & MHRA & No marketing authorization; trial and device documentation aligned to MHRA framework \\
Canada & Health Canada & No marketing authorization; dossier structured for future Health Canada extension \\
Australia & TGA & No marketing authorization; safety and audit-trail design aligned to TGA requirements \\
\bottomrule
\end{tabularx}
\end{table}

\figcaption{Table 2. Cross-jurisdiction status of daraxonrasib and the LLM-directed
device. The agent and device are investigational in every jurisdiction; the
documentation is authored once to align with the principal international device and
trial frameworks so it can support future multi-jurisdiction extension.}

\subsection{Overview of Preclinical Data}

The preclinical and computational basis for this investigation supports both the
pharmacologic rationale for daraxonrasib in KRAS G12-mutated PDAC and the safety and
feasibility rationale for the LLM-directed robotic device, and is presented in full
in the Pharmacology and Toxicology information (\S 8) and the Investigator's
Brochure (\S 5). On the pharmacology side, daraxonrasib acts through the
mechanism-based, on-target suppression of RAS pathway signaling described in
\S 3.1.2: by engaging the active GTP-bound state of mutant RAS in a tri-complex with
cyclophilin A, it suppresses MAPK and PI3K effector signaling and reduces tumor-cell
proliferation and survival across the KRAS G12 alleles that define the trial
population \cite{daraxwhipple}. The drug-distribution, exposure, and tolerability
assumptions that anchor the three dose levels (DL1 160 mg, DL2 220 mg, DL3 300 mg)
are carried forward from the pan-KRAS RASolute 302 program, in which 300 mg once
daily is the recommended Phase 2 dose, and the anticipated mechanism-based toxicity
profile (gastrointestinal, dermatologic and mucocutaneous, and hepatic) informs the
DLT definitions and the perioperative pause-and-restart logic \cite{onpremwhippl}.

On the device and combined-intervention side, the preclinical evidence is dominated
by a structured program of high-fidelity simulation and digital-twin work that
substitutes for and supplements conventional preclinical testing for the Physical AI
system, consistent with the verification, validation, and uncertainty
quantification (VVUQ) discipline that the protocol applies. The mandatory Phase 0
simulation-validation gate requires at least 1000 simulated procedures across at
least two independent frameworks (for example Isaac, MuJoCo, Gazebo, and PyBullet),
a Unified Safety Level rating of at least 7.0, and a sim-to-real gap below 2 mm in
trajectory and below 0.5 N in tip force, with re-validation after any qualifying
change; this gate is the principal preclinical safety bar that the device must clear
before any patient-facing activity. The supporting computational evidence base
includes the 60-second PDAC robotic Whipple and daraxonrasib simulation, the AI
digital-twin pancreatic-cancer simulation for in silico FDA compliance and cost
efficiency \cite{fdadigtwinpc}, the QSP metastatic-PDAC trial simulation with code,
VVUQ, and a reproducible playbook \cite{qspmetpancre}, the 100,000-patient in silico
Phase 3 five-arm PDAC trial triplicate \cite{chatgpt100kp}, and the end-to-end PDAC
digital-twin clinical-trial proposals \cite{pdacdigtwinp}. Together these establish
that the integrated drug-device intervention has been characterized in silico across
a large and reproducible body of simulated experience before any human exposure, and
that the on-premises advisory LLM operates strictly as a bounded, hash-pinned,
second-opinion oracle with full autonomy prohibited consistent with 21 CFR
\S 312.21(e) \cite{clintrustpai, congress9510}. The integrated pharmacology and
toxicology summary, including the structure of the pharmacology summary, the
integrated toxicology summary, and the full data tabulation, is presented in \S 8.

\subsection{References}

The references cited throughout this Introduction are drawn from the IND reference
library and are collected, formatted, and rendered in the References and Back Matter
section of this dossier. The principal sources grounding the clinical subject, the
daraxonrasib mechanism and class, the device safety architecture, the indication,
the duration timeline, and the acceleration thesis include the FDA Phase 1 content
and format guidance \cite{phase1guidance}, the oncology-trial LLM authoring basis
\cite{oncotrialllm}, the daraxonrasib robotic Whipple and on-premises device sources
\cite{daraxwhipple, onpremwhippl}, the 21 CFR Part 312, Part 50, and ICH Physical AI
unification adaptations \cite{cfr312paiuni, cfr050paiuni, ichpaistduni}, the PDAC
epidemiology and acceleration sources \cite{pdac060s2030, paipredict4x,
natlpaiplatf}, the verification-before-generation and clinical-trust Physical AI
sources \cite{congress9510, clintrustpai, paiautospons, paisitedocpk}, the
comparative LLM evaluation against the FDA Adverse Event Reporting System
\cite{llmcomp16fda}, and the computational digital-twin and QSP evidence base
\cite{fdadigtwinpc, qspmetpancre, chatgpt100kp, pdacdigtwinp}. Each citation key
resolves against the IND reference library, and the assembled list renders in the
References and Back Matter section.
